\documentclass{ximera}
% xmPreamble.tex en xmPrintstyle.sty worden automatisch geladen door ximera.cls (v2.7.8);
% NIET nog eens \input{./preamble.tex} of \addPrintStyle{.} doen (geeft 'already defined'-fouten).

\makeatletter % eigen img-map toevoegen aan het zoekpad (werkt standalone én vanuit de xourse)
\g@addto@macro\Ginput@path{{img/}{UW4202/img/}}
\makeatother

\begin{document}
\author{Alexander Holvoet}
\xmtitle{Een hoop zand op een ronde ondergrond}{}

Giet zo veel mogelijk zand of zout op een ronde basisvorm, bijvoorbeeld een rond
bierviltje of een omgekeerd plastieken bekertje. Je moet het niet al te voorzichtig doen:
wat te veel is, glijdt er af. Daarom staat je basisvorm best een beetje hoger.

\begin{exercise}
    Welke vorm hebben de hopen?
    \begin{oplossing}
        Het zijn allemaal kegels. Je zult merken dat deze kegels niet noodzakelijk perfect
        zijn: de top is niet juist \'e\'en punt. We werken hier met een \textit{wiskundig
        model} dat vrij goed aansluit bij de realiteit, maar er niet helemaal mee samenvalt.
    \end{oplossing}
\end{exercise}

We definiëren een nieuw begrip: de \emph{rusthoek} is de hoek tussen de beschrijvende
rechten van de kegel en de horizontale richting. Die hoek is over de hele flank gelijk,
dus niet enkel vanonder.

% TODO (afbeelding ontbreekt): in de bron stond hier \afbeelding[5cm]{Lrusthoek} ter
% illustratie van de rusthoek. Deze lijntekening zit niet (meer) in img/. Voeg ze toe
% wanneer beschikbaar.

\begin{exercise}
    Denk even na hoe je de rusthoek zo nauwkeurig mogelijk kunt meten.
    \begin{hint}
        Een geodriehoek lijkt voor de hand te liggen, maar zodra je de hoop raakt, begint
        die te schuiven. Kun je de hoek meten \emph{zonder} de zandhoop aan te raken?
    \end{hint}
    \begin{oplossing}
        Een geodriehoek gebruiken is lastig: zodra je de hoop raakt, begint die te schuiven.
        Beter neem je een foto van de kegel in vooraanzicht en meet je daarop de rusthoek,
        of je plakt de foto in GeoGebra en meet de rusthoek daar. Je kunt ook de straal
        van het grondvlak en de hoogte van de kegel meten en daarmee de tangens van de
        rusthoek \textemdash{} en dus de rusthoek zelf \textemdash{} berekenen.
    \end{oplossing}
\end{exercise}

% TODO (afbeelding ontbreekt): in de bron stond hier
% \afbeelding[6cm]{rusthoek meten met geogebra}. Niet aanwezig in img/.

\begin{exercise}
    Meet en noteer de rusthoeken voor hopen met dezelfde diameter maar met verschillend
    materiaal. Verdeel eventueel het werk. Wat besluit je?
    \begin{oplossing}
        Een voorbeeld van meetresultaten:
        \begin{center}
        \begin{tabular}{l l l}
            \textbf{diameter} & \textbf{materiaal} & \textbf{rusthoek} \\
            \(10\) cm & zand  & \(32^\circ\) \\
            \(10\) cm & rijst & \(25^\circ\) \\
            \(10\) cm & zout  & \(30^\circ\) \\
        \end{tabular}
        \end{center}
        Als je de diameter van het grondvlak constant houdt, dan hangt de rusthoek af van
        het soort materiaal.
        \begin{image}[6cm]
            \includegraphics[width=6cm,keepaspectratio]{lmateriaal.jpg}
        \end{image}
        Verschillende materialen geven een verschillende rusthoek.
    \end{oplossing}
\end{exercise}

\begin{exercise}
    Doe nu hetzelfde voor hopen van hetzelfde materiaal maar met verschillende diameters.
    Meet en noteer ook hier je meetresultaten in een tabel. Wat besluit je?
    \begin{oplossing}
        Een voorbeeld van meetresultaten:
        \begin{center}
        \begin{tabular}{l l l}
            \textbf{materiaal} & \textbf{diameter} & \textbf{rusthoek} \\
            zand & \(10\) cm & \(32^\circ\) \\
            zand & \(15\) cm & \(32^\circ\) \\
            zand & \(20\) cm & \(32^\circ\) \\
        \end{tabular}
        \end{center}
        De rusthoek hangt niet af van de diameter van het grondvlak.
    \end{oplossing}
\end{exercise}

\end{document}
